\chapter{App A --- EVE-NG Installation on Proxmox VE}

This appendix documents the configuration used to deploy EVE-NG Community Edition on the Proxmox VE server used for this project. The procedure has been validated and the same steps can be used to reproduce the environment on any compatible Proxmox host. Screenshots and extended notes are available in the companion web documentation at docs/web/docs/guides/eve\_ng\_install\_proxmox/.

\section{Virtual Machine Requirements}

\begin{longtable}[]{@{}
  >{\raggedright\arraybackslash}p{(\columnwidth - 2\tabcolsep) * \real{0.5000}}
  >{\raggedright\arraybackslash}p{(\columnwidth - 2\tabcolsep) * \real{0.5000}}@{}}
\toprule\noalign{}
\begin{minipage}[b]{\linewidth}\raggedright
\textbf{Parameter}
\end{minipage} & \begin{minipage}[b]{\linewidth}\raggedright
\textbf{Value used in this project}
\end{minipage} \\
\midrule\noalign{}
\endhead
\bottomrule\noalign{}
\endlastfoot
Hypervisor & Proxmox VE 8.x \\
CPU type & \texttt{host} (required for nested virtualisation) \\
vCPU cores & 16 (minimum 8) \\
RAM & 32\,GB per EVE-NG instance (host upgraded to 64\,GB total) \\
Disk & 100\,GB VirtIO, SSD-backed Proxmox pool \\
SCSI controller & VirtIO SCSI \\
Network & VirtIO, bridged on \texttt{vmbr0} (management); \texttt{pnet0}--\texttt{pnetN} as cloud bridges \\
BIOS & SeaBIOS (default; OVMF also works) \\
\end{longtable}

\section{Installation Steps}

\begin{itemize}
\tightlist
\item
  Download the EVE-NG Community Edition ISO from the official website.

  \begin{itemize}
  \tightlist
  \item
    In Proxmox, create a new VM with the parameters in Table . Ensure the CPU type is set to \texttt{host} so that virtualisation extensions (VT-x / AMD-V) are exposed to the guest.
  \item
    Boot from the ISO and follow the Ubuntu Server installation wizard. Assign a static IP address and install the OpenSSH server during setup.
  \item
    After the first reboot, follow the EVE-NG Community post-install wizard (\texttt{/opt/unetlab/wrappers/unl\_wrapper\ -a\ postinstall}) and configure the management interface.
  \item
    Run \texttt{apt\ update\ \&\&\ apt\ upgrade\ -y} and reboot once more to ensure all EVE-NG patches are applied.
  \item
    Access the web interface at \texttt{https://\textless{}EVE-NG-IP\textgreater{}/} and change the default credentials immediately.
  \end{itemize}
\end{itemize}

\section{Proxmox Network Configuration}

Two bridge types are used:

\begin{itemize}
\tightlist
\item
  **\texttt{vmbr0**\ (bridged)} --- management access from the LAN; also used as \texttt{pnet0} inside EVE-NG for labs that require connectivity to the homelab network (e.g., Lab 1).

  \begin{itemize}
  \tightlist
  \item
    **\texttt{vmbr1**\ (NAT)} --- optional internet-facing bridge to keep lab traffic isolated from the LAN while allowing outbound connectivity.
  \end{itemize}
\end{itemize}

\section{Post-Install Checklist}

\begin{itemize}
\tightlist
\item
  \textbullet{} CPU type set to \texttt{host} in Proxmox VM options.

  \begin{itemize}
  \tightlist
  \item
    \textbullet{} VM autostart enabled in Proxmox.
  \item
    \textbullet{} Snapshot taken after successful EVE-NG installation.
  \item
    \textbullet{} Default admin password changed.
  \item
    \textbullet{} OS Client Side installed on instructor workstation (HTML5 console option does not require it).
  \item
    \textbullet{} Image directories created under /opt/unetlab/addons/qemu/ and permissions fixed with \texttt{unl\_wrapper\ -a\ fixpermissions}.
  \end{itemize}
\end{itemize}
